\section{Introduction}
Modern large language model (LLM)-based agents interact with external environments through iterative cycles of reasoning, action, and observation~\citep{yao2023react}. They increasingly rely on stateful tools that inspect or modify evolving runtime environments, extending beyond simple
information retrieval. 
Frameworks such as Claude Code and OpenClaw decouple
model serving from the tool-execution environment
~\citep{anthropic2026claudecode,openclaw2026runtime}, forming what we term a disaggregated agent loop. Because the standard loop serializes model generation
and tool execution, tools account for a growing share of end-to-end latency as decoding accelerates~\citep{tilert2026, artificialanalysis_anthropic, artificialanalysis_openai, groq2024specdec, cerebras2025k2think}.

Speculation can break this serialization by starting likely future work early.
Existing systems target one side of the loop: \textit{(i)} action speculation
predicts and launches future tool calls
\citep{ye2025speculative,sui2026parallelizing,spectoolcalls2025,bpaste2026,spork2026},
whereas \textit{(ii)} observation speculation continues generation from
provisional observations verified asynchronously~\citep{saberi2026spechop}.

However, existing speculation methods fall short on heterogeneous stateful tool workloads, exposing three barriers to effective speculation.
\textbf{(1) Latency concentration:} a few slow calls dominate tool time, decoupling hit rate from latency savings (Section~\ref{sec:challenge1}).
\textbf{(2) Environment dependence:} some environment-dependent observations emerge only through execution, making observation prediction alone insufficient and requiring sandboxed action speculation (Section~\ref{sec:challenge2}).
\textbf{(3) Lookahead tradeoff:} launching speculative actions earlier increases their overlap with model generation, but lossless generation requires the intervening action chain to match;
compounding errors collapse accuracy and erase the potential speedup
(Section~\ref{sec:challenge3}). 
Together, these findings demand a new approach
to speculation across the full agent--environment loop.

We present \sysname{}, a lossless framework that jointly speculates actions and observations across the disaggregated agent loop
(Section~\ref{sec:method}). 
To address latency concentration, \emph{Expected
Value Decoding} (\evd) generates observation candidates by combining their probability and estimated tool time, directly optimizing expected time hidden rather than hit rate. 
For observations only execution can reveal, \sysname{} launches latency-critical target actions in isolated forks that contain their effects~\citep{cubesandbox_2026,deltabox2026}. 
To break the lookahead--accuracy tradeoff, \emph{Joint Action--State Verification} (JASV) recasts long-horizon action dependency from full-chain prediction into target action--state verification: a fork is verified and selected only when its predicted action
and origin state match the emitted action and committed state. This unlocks
long-range overlap without predicting intervening actions, while exact
verification preserves serial semantics and lossless. Together, these mechanisms yield
11.8--32.5\% end-to-end latency savings, with gains increasing as decoding
accelerates as shown in Figure~\ref{fig:tpot}.

% \newpage
We make three contributions:
\begin{itemize}
\item We characterize and quantify the latency concentration, information
boundary, and lookahead tension that limit speculation over stateful tools.
\item We introduce \sysname{}, combining \evd{} for latency-aware observation
drafting, isolated action execution for environment-dependent observations, and
JASV for lossless long-range lookahead without full-chain prediction.
\item Across 4 different harnesses and 5 models on Terminal-Bench, \sysname{} achieves 11.8--32.5\% end-to-end latency savings, reduces p99 latency by up to 42.8\% and outperforms practical baselines.
\end{itemize}
